\chapter{定义、模型与理论基础}

\section{BCI-IoT 通信对象与安全资产}
本文研究对象为位于应用层与 TCP 之间的传输增强层。该层面向 BCI-IoT 场景中的连续脑电数据上传、控制指令回传和反馈参数同步，目标是在不改动内核 TCP 的条件下提供会话保护、异常处置和外观调控能力。系统需支持上行脑电片段或类别结果与下行反馈指令并发存在的业务模式。

本文将 BCI-IoT 通信资产划分为业务数据、会话材料、协议状态和外观特征四类，如表 \ref{tab:c2-assets} 所示。其中，外观特征被列为独立资产，是因为脑电任务类别、控制意图和反馈频率可能映射到报文长度、到达间隔和字节统计特征。已有研究表明，BCI 信号本身可成为隐私推断和侧信道攻击对象\cite{martinovic2012sidechannel,lange2018pinbci}；在联网场景中，通信外观会进一步扩大此类风险的可观测范围。

\begin{table}[H]
\centering
\small
\caption{BCI-IoT 通信资产与安全属性}
\label{tab:c2-assets}
\begin{tabular}{p{2.5cm}p{5.0cm}p{5.2cm}}
\toprule
资产类型 & 主要内容 & 目标属性 \\
\midrule
业务数据 & 脑电类别、控制参数、反馈指令、用户状态 & 机密性、完整性、语义最小暴露 \\
\midrule
会话材料 & 临时密钥、派生密钥、nonce salt、确认密钥 & 密钥新鲜性、密钥隔离、会话一致性 \\
\midrule
协议状态 & 握手摘要、认证结果、重放缓存、异常记录 & 状态一致性、重放可检测、异常可终止 \\
\midrule
外观特征 & 包长、时序、方向、可打印比例、字节熵 & 类别弱相关性、协议识别难度、统计可调性 \\
\bottomrule
\end{tabular}
\end{table}

\section{威胁模型与研究范围}
\subsection{系统模型}
本文采用“采集端--网络链路--服务端”三段式系统模型。采集端负责产生 BCI 数据或控制请求，并通过 TCP 连接接入服务端；服务端负责完成认证握手、会话密钥建立、记录层解封装与业务处理；网络链路可能经过无线接入、NAT、防火墙、公网中继或代理节点。通信双方均保存必要的会话状态，协议状态规模需满足资源受限设备的部署要求。

本文讨论的通信增强层包含三类机制：认证握手机制、AEAD 记录层和外观编码层。认证握手机制提供身份认证、临时密钥协商、transcript 绑定与重放检测；AEAD 记录层提供业务明文的机密性与完整性；外观编码层通过组合数独提示字节和随机填充改变可观测字节分布与长度特征。

\subsection{攻击者能力}
本文采用“被动观测 + 主动干扰”的组合威胁模型。被动攻击者可长期抓取流量并训练分类模型，观测内容包括连接元数据、包长、方向、到达间隔、可打印比例和字节熵等外部特征。主动攻击者可重放历史握手或会话报文，篡改握手字段，注入畸形探测流量，或尝试通过协议错误行为区分真实服务与异常处置路径。在标准密码算法安全成立的前提下，本文重点分析协议工程层面的可识别性、重放风险、异常路径和资源消耗。

\begin{table}[H]
\centering
\small
\caption{本文关注的主要威胁}
\label{tab:c2-threat-model}
\begin{tabular}{p{3.1cm}p{5.3cm}p{5.0cm}}
\toprule
威胁类型 & 攻击能力 & 主要影响 \\
\midrule
重放攻击 & 记录并重复发送历史握手或会话报文 & 触发重复控制，扰乱会话状态 \\
MITM 篡改 & 插入链路并修改握手或业务字节 & 身份冒充、数据篡改、会话破坏 \\
被动侧写 & 基于长度、时序、字节分布建模分类 & 协议识别与业务类别推断 \\
DPI 探测 & 基于可打印率与规则进行判别 & 协议暴露与策略阻断 \\
资源耗尽 & 持续发送畸形探测流量 & CPU/内存资源异常增长 \\
\bottomrule
\end{tabular}
\end{table}

\subsection{研究假设与范围}
为保证论证可检验，本文明确以下假设：
\begin{itemize}
  \item \textbf{A1（密码原语安全）}：不讨论标准密码算法被直接攻破情形，重点分析协议工程与实现路径风险；
  \item \textbf{A2（端侧资源受限）}：设备侧内存与 CPU 预算有限，协议状态必须有上界；
  \item \textbf{A3（公网受限可达性）}：真实部署中可能存在 UDP 受限、代理终止与策略审查\cite{rfc8085,rfc4787,nedergaard2023coapsecure}；
  \item \textbf{A4（可观测外观）}：攻击者可长期获得长度、时序、方向与字节分布等统计特征。
\end{itemize}

表 \ref{tab:c2-risk-scope} 给出本文对 BCI-IoT 风险的处理范围。通信层能够处理链路窃听、握手篡改、重放、探测和外观侧写；采集硬件失陷、终端系统被完全控制、脑电模型投毒和医疗合规治理属于其他层面的研究任务，本文在第五章按限制条件讨论。

\begin{table}[H]
\centering
\small
\caption{BCI-IoT 风险与本文处理范围}
\label{tab:c2-risk-scope}
\begin{tabular}{p{3.2cm}p{4.7cm}p{5.1cm}}
\toprule
风险来源 & 本文处理方式 & 处理原因 \\
\midrule
脑电数据泄露 & AEAD 记录层保护业务明文 & 脑电类别和控制参数具有隐私属性 \\
握手篡改与身份冒充 & 证书签名、摘要绑定和确认码 & 闭环控制需要确认通信对象与会话密钥一致 \\
历史报文重放 & 时间窗口、随机数和重放缓存 & 重复反馈或重复控制会扰乱会话状态 \\
加密流量侧写 & 外观编码、随机填充与类别恢复实验 & 长度和时序可能暴露脑电任务类别 \\
端侧资源消耗 & 有界状态、异常路径和内存指标 & 可穿戴与边缘设备资源受限 \\
\bottomrule
\end{tabular}
\end{table}

\section{基础密码机制与会话保护}
认证加密（Authenticated Encryption with Associated Data, AEAD）是一种同时提供数据机密性、完整性和真实性的加密模式\cite{rfc5116}。形式化地，一个 AEAD 方案由三个算法组成：密钥生成算法 $\mathcal{K}$、加密算法 $\mathcal{E}$ 和解密算法 $\mathcal{D}$。加密过程可表示为
\[
C = \mathcal{E}_K(N, A, P),
\]
其中 $K$ 为密钥，$N$ 为随机数（Nonce），$A$ 为关联数据（Associated Data），$P$ 为明文。解密过程表示为
\[
P/\bot = \mathcal{D}_K(N, A, C),
\]
若验证失败则输出符号 $\bot$。

AEAD 接口能够保护单条记录的机密性与完整性，但身份认证、密钥新鲜性、重放检测和异常终止条件需要由握手状态机提供。TLS/DTLS 的握手设计表明，安全通道建立通常需要身份认证、临时密钥协商、transcript 绑定和记录层密钥派生共同完成\cite{rfc8446,rfc9147}。本文协议设计参考上述标准化实践，在轻量原型中实现证书签名、临时密钥协商、摘要绑定和确认码，并将其与组合数独外观层组合。

\section{流量侧写与外观约束}
流量分析（Traffic Analysis）是指攻击者在无法解密加密流量内容的情况下，通过观察数据包的外部特征推断通信内容、用户身份或应用类型的攻击技术\cite{alwhbi2024encrypted,hou2024encmalware}。根据攻击目标的不同，流量分析可分为协议识别、网站指纹识别和用户行为分析。加密机制是保障通信机密性与隐私的重要基础，但也会削弱传统基于载荷内容的流量检查能力\cite{fu2025encryptedtraffic,hou2024encmalware}。

在加密流量中，虽然载荷内容被隐藏，但协议交互产生的侧信道特征往往难以完全消除：
\begin{itemize}
  \item \textbf{包长特征（Packet Size）}：不同应用层协议或业务逻辑产生的数据包大小分布具有显著差异；
  \item \textbf{时序特征（Inter-arrival Time, IAT）}：数据包到达的时间间隔反映应用层的交互频率与处理延迟；
  \item \textbf{方向特征（Direction）}：上下行流量的比例与突发模式构成流量方向性指纹；
  \item \textbf{字节统计特征}：字节熵、可打印比例和固定魔数可用于 DPI 规则或协议识别。
\end{itemize}

针对 BCI-IoT 场景，攻击者可利用这些特征推断用户当前的脑电状态或控制指令类别。本文采用与加密流量分类研究相近的评估流程：先抽取长度、字节熵、可打印比例和读写调用统计，再执行协议识别与类别恢复评估。TLS 中间人代理解密对照用于刻画特定运维/攻防条件下终止代理点的可见性风险，即“链路加密已启用但业务仍可恢复”的现实情况\cite{wilkens2021passivetls}。

\section{纯组合数独数学基础与术语}
\subsection{术语定义}
为减少歧义，本文统一采用以下中文术语：
\begin{itemize}
  \item \textbf{提示字节}：编码后可被解码器识别为“数独提示位”的字节；
  \item \textbf{填充字节}：仅用于外观扰动、解码时被忽略的字节；
  \item \textbf{模式模板}：定义字节位上“冗余位/位置位/取值位”分布的 8 位模式；
  \item \textbf{解码键}：4 个提示字节按序拼接得到的 32 位键值。
\end{itemize}

在实现与论文叙述中，本文将“提示字节/填充字节/模式模板”界定为功能术语。安全性论证建立在标准握手、密钥派生、AEAD 与状态机正确性之上，外观编码承担外观调控功能，密码学安全仍由握手、密钥派生和记录层提供。

\begin{table}[htbp]
\centering
\small
\caption{纯组合数独外观层符号}
\label{tab:c2-symbols}
\begin{tabular}{p{2.3cm}p{8.9cm}}
\toprule
符号 & 含义 \\
\midrule
$\mathcal{B}$ & 字节集合，$\mathcal{B}=\{0,\ldots,255\}$ \\
$\mathcal{S}$ & 满足 4$\times$4 组合数独约束的解空间 \\
$f_k$ & 会话密钥 $k$ 下映射，$f_k:\mathcal{B}\rightarrow\mathcal{S}$ \\
$\mathcal{C}$ & 提示位组合集合，$|\mathcal{C}|=\binom{16}{4}$ \\
$D$ & 解码键到原字节的映射表 \\
$M$ & 填充字节池大小 \\
$n$ & 明文字节长度 \\
\bottomrule
\end{tabular}
\end{table}

4$\times$4 组合数独解空间定义为
\[
\mathcal{S}=\{g\mid g \text{满足行、列与 }2\times2\text{ 子宫格约束}\},
\]
该集合规模为 $|\mathcal{S}|=288$，可选取其中 256 个解与字节集合建立一一对应\cite{oeisA107739,oeisA109252}。

\subsection{组合空间与排列基数}
从组合数学角度看，若按会话种子从 288 个可行网格中为 256 个字节选择互异映射，则可行映射上界为排列数
\[
P(288,256)=\frac{288!}{32!}\approx 2^{1825}.
\]
该值用于刻画外观映射重排的数量级上界，可反映离散映射空间规模。进一步地，每个字节的提示位位置组合满足
\[
|\mathcal{C}|=\binom{16}{4}=1820,
\]
与填充注入共同作用时，会显著扩大可观测外观形态数量。该分析与数独相关组合研究关于搜索空间庞大且约束主导求解复杂度的结论一致\cite{lynce2006satgrid}。

\subsection{可逆映射与复杂度}
对字节 $b\in\mathcal{B}$，先映射到网格 $g_b=f_k(b)$，再从网格中抽取 4 个提示位编码为 4 个提示字节。解码时按 4 字节拼接形成解码键，通过映射表 $D$ 还原原字节。构表阶段若发现不同原字节产生同一解码键，则直接拒绝构表，因此可保证单值可逆。

在 ASCII 布局下，提示字节可写为
\[
h = 0x40 \,|\, (v \ll 4)\,|\, p,
\]
其中 $p\in[0,15]$ 为位置编码，$v\in[0,3]$ 为取值编码。编码与解码对长度为 $n$ 的消息均为线性复杂度 $O(n)$。

无填充时，若第 $i$ 个字节候选提示组数量为 $K_i$，可观测形态规模为
\[
N_{\text{无填充}}=\prod_{i=1}^{n}K_i\approx 8^n.
\]
若每字节存在 5 个填充注入检查点，报文尾部 1 个检查点，填充池大小为 $M$，则
\[
N_{\text{线上}} = N_{\text{无填充}}(M+1)^{5n+1}.
\]
该表达式解释了“外观随机性增强”与“带宽开销上升”之间的必然权衡。

\section{统一评测指标与统计规则}
本文采用如下统一指标：
\[
R_{\text{开销}}=\frac{B_{\text{wire}}}{B_{\text{payload}}},
\quad
\overline{\mathrm{RTT}}=\frac{1}{N}\sum_{i=1}^{N}\mathrm{RTT}_i,
\quad
\mathrm{P95}=\operatorname{Quantile}_{0.95}(\{\mathrm{RTT}_i\}),
\]
\[
H=-\sum_{x=0}^{255}p_x\log_2 p_x,
\quad
\rho_{\text{可打印}}=\frac{\sum_{x=0x20}^{0x7E}c_x}{\sum_{x=0}^{255}c_x},
\quad
T_{\text{payload}}=\frac{B_{\text{payload}}}{t}.
\]
其中 $H$ 为字节熵，$\rho_{\text{可打印}}$ 为可打印字节比例。所有关键结论均要求“多次运行 + 中位数 + 稳定性检查”，避免单次抖动主导结果。

\section{本章小结}
本章给出了后续协议设计与实验评估所需的定义、模型与约束：首先明确 BCI-IoT 通信资产和安全属性，其次给出系统模型、威胁模型和研究范围，随后说明 AEAD、流量侧写、组合数独外观编码与统一评测指标。这些内容构成第 3 章实现和第 4 章实验分析的基础。
